\chapter{Composite inputs and finite cubic models}
\label{chap:composite-operations}

A higher operation on generators can be one Taylor coefficient.
An operation on their products can multiply several such coefficients.
This distinction determines how much of a potential is retained by
its boundary algebra. It also determines when the complete family
of higher operations is finite.

The ordered diagonal of Chapter~\ref{chap:ordered-taylor} gives a
concrete setting for this question. Its comparison is defined on
every exterior monomial, not only on the primitive generators.
The first composite calculation already differs from the primitive
formula. A sum over explicit contractions gives every composite
value. Counting those contractions then bounds the possible arities.

\section{A product of two Taylor coefficients}
\label{sec:co-composite-example}

Use the ordered right-action diagonal with two variables. Thus
\(R=k[y_1,y_2]\), \(S=R[r_1,r_2]\), \(x_i=y_i+r_i\), and
\(k\) is a commutative rational algebra. The formal version completes
in the displayed variables. The odd endomorphisms \(H_i,\beta_i\)
and the central coefficients \(c_{ij}\) satisfy
\[
 \{H_i,H_j\}=c_{ij},\qquad
 \{H_i,\beta_j\}=\delta_{ij},\qquad
 \{\beta_i,\beta_j\}=0.
\]
Here \(\{A,B\}=AB+BA\) on odd endomorphisms. The ordered monomials
\(H_I\beta_J\) form the basis used by the contraction
\eqref{eq:ot-homotopy}. Its projection \(p\) sets \(r=0\), kills
each term with a \(\beta\), and sends \(H_I\) to \(\xi_I\).
Write \(e=\mathsf{s}1\), \(a_i=\mathsf{s}\xi_i\), and
\(\chi=\mathsf{s}(\xi_1\xi_2)\), where \(\mathsf{s}\) reverses parity.
The \(a_i\) are even; \(e\) and \(\chi\) are odd.

For the divided difference \(\Delta_j\) in \eqref{eq:ot-homotopy},
put \(W_{(a,b)}=\partial_1^a\partial_2^bW(y)/(a!b!)\).
The first formula with two composite inputs is
\begin{equation}\label{eq:co-first-composite}
 b_3(\chi,\chi,a_1)=-W_{(3,0)}W_{(0,2)}e.
\end{equation}
It holds for every polynomial or formal potential, with its
retained parameters unchanged.

Here is the calculation. Put \(K=H_1H_2\). The Clifford relations give
\[
 K^2=-\tfrac14c_{11}c_{22}+c_{12}K.
\]
Since \(K\) is even, the second suspended comparison component is
\[
 F_2(\chi,\chi)=\mathsf{s}\left(
 \tfrac14\sum_j\Delta_j(c_{11}c_{22})\beta_j
 -\sum_j\Delta_j(c_{12})K\beta_j\right).
\]
Every higher component takes values in the left ideal
\(\mathcal J=\sum_jE\beta_j\). Indeed, the homotopy adds a
\(\beta\) on the right in the ordered basis.
Multiplying an element of \(\mathcal J\) on the left keeps it
in \(\mathcal J\), and \(p\mathcal J=0\).
Thus only the last-letter split can contribute to \(b_3\).
Its sign is negative, because the matrix part of \(F_2(\chi,\chi)\)
is odd. In multiplying by \(H_1\), only
\(\beta_1H_1=1-H_1\beta_1\) can give a scalar.

For this order, \(c_{12}\) and \(c_{22}\) are independent of \(x_1\).
Furthermore
\[
 (\Delta_1c_{11})|_{r=0}=-2W_{(3,0)},\qquad
 c_{22}|_{r=0}=-2W_{(0,2)}.
\]
These identities follow by differentiating the one-variable divided
difference defining \(g_1\), and by evaluating \(c_{22}\) on the
diagonal. Consequently the projected scalar is
\(-\tfrac14(-2W_{(3,0)})(-2W_{(0,2)})\), proving
\eqref{eq:co-first-composite}. Two scalar contractions contribute.
The product of their coefficients is part of the operation.

\section{Multiplication as a finite exterior differential operator}
\label{sec:co-symbol-product}

To compute arbitrary products, replace the ordered matrix basis by
exterior symbols. For any number \(d\) of variables put
\[
 \mathcal P=S\otimes_k\Lambda(\theta_1,\ldots,\theta_d,
                                      \epsilon_1,\ldots,\epsilon_d).
\]
The symbol \(\theta_I\epsilon_J\) represents \(H_I\beta_J\).
Exterior multiplication is denoted by \(\mu\). It is initially
different from multiplication of the represented endomorphisms.

List the symbols as \(q_1,\ldots,q_{2d}\), with the \(\theta\)'s first.
Their left odd derivatives \(\partial_i\) delete \(q_i\), with the
sign of its position, and give zero if it is absent. Define the
central matrix \(\Gamma\) by
\[
 \Gamma_{ij}=
 \begin{cases}
 c_{ij},&1\le j<i\le d,\\
 c_{ii}/2,&i=j\le d,\\
 \delta_{ab},&i=d+a,\ j=b\le d,\\
 0,&\text{otherwise}.
 \end{cases}
\]
The entry \(\Gamma_{ij}\) is the scalar contraction when the left
symbol \(q_i\) crosses the right symbol \(q_j\).
For example, \(\Gamma_{d+i,i}=1\) records
\(\beta_iH_i=1-H_i\beta_i\).

Tensor products of odd maps use the Koszul sign:
\((\partial_i\otimes\partial_j)(f\otimes g)
 =(-1)^{|f|}\partial_i f\otimes\partial_jg\) for homogeneous \(f\).
Set
\begin{equation}\label{eq:co-symbol-star}
 f\star g=\mu\exp\left(-\sum_{i,j=1}^{2d}
                   \Gamma_{ij}\partial_i\otimes\partial_j\right)
                   (f\otimes g).
\end{equation}
The exponential is finite: each application deletes one symbol
from each input. It is therefore defined over every rational
coefficient algebra, without convergence assumptions.

\begin{proposition}\label{prop:co-symbol-product}
The symbol map \(\theta_I\epsilon_J\mapsto H_I\beta_J\) is an
\(S\)-algebra isomorphism from \((\mathcal P,\star)\) to \(E\).
For homogeneous \(f\), its product has the explicit expansion
\[
 f\star g=
 \sum_{k=0}^{2d}\frac{(-1)^{k(k+1)/2+k|f|}}{k!}
 \sum_{\substack{i_1,\ldots,i_k\\j_1,\ldots,j_k}}
 \left(\prod_{a=1}^k\Gamma_{i_aj_a}\right)
 (\partial_{i_1}\cdots\partial_{i_k}f)
 (\partial_{j_1}\cdots\partial_{j_k}g),
\]
where every index ranges from \(1\) to \(2d\), and the final
product is exterior multiplication.
\end{proposition}

\begin{proof}
The sign in the expansion has three contributions. The minus
in the exponential contributes \((-1)^k\). Passing right-factor
derivatives across subsequent left-factor derivatives contributes
\((-1)^{k(k-1)/2}\). Acting on \(f\otimes g\) contributes
\((-1)^{k|f|}\). Their product is the stated sign.

For three factors, let \(T_{ab}\) be the exponent's operator
acting on factors \(a,b\). The operators \(T_{12},T_{13},T_{23}\)
commute: each has even degree in the anticommuting odd derivatives,
and no derivative acts on a coefficient \(\Gamma_{ij}\).
The odd product rule shows that both parenthesizations of three
products are exterior multiplication after
\(\exp(T_{12}+T_{13}+T_{23})\). Thus \(\star\) is associative.
The constant symbol \(1\) is its unit.

On generators, \(q_i\star q_j=q_iq_j+\Gamma_{ij}\).
These equations give exactly the stated Clifford relations.
An increasing word in the \(q_i\) acquires no contraction when
multiplied from left to right. Hence its \(\star\)-product is
its exterior monomial. The resulting map to \(E\) sends an
exterior basis to the already established basis \(H_I\beta_J\).
It is therefore an algebra isomorphism.
\end{proof}

Antisymmetrized representatives have an equally explicit formula.
Replace the upper-left block of \(\Gamma\) by \(c_{ij}/2\) for
all \(i,j\le d\), keeping the other blocks unchanged.
Then \(\theta_I\) represents the antisymmetrized product \(H_{[I]}\).
To verify this assertion, antisymmetrize the product of distinct
\(\theta\)'s in \eqref{eq:co-symbol-star}. Every term containing
a contraction cancels, since its contracted coefficient is symmetric
in those two indices. The remaining term is the exterior monomial.
The same associativity proof applies. For antisymmetrized representatives,
the sign of \(c_{ij}\) is still fixed by their actual anticommutator,
and the normal-coordinate contraction retains its chosen order.

\section{A scalar formula on every input}
\label{sec:co-tree-formula}

In exterior symbols, \(p\) kills every nonempty \(\epsilon\)-monomial
and evaluates coefficients at \(r=0\). For the increasing-order
contraction, the homotopy is
\[
 h(f\theta_I\epsilon_J)=(-1)^{|I|}
       \sum_{j<\min J}\Delta_j f\,\theta_I\epsilon_j\epsilon_J,
 \qquad \min\varnothing=d+1.
\]
Thus multiplication, projection, and homotopy are explicit scalar
operations on polynomials or formal series. No matrix inversion is needed.

A planar binary tree is a rooted tree whose internal vertices each
have an ordered left and right input. Its leaves are read from left
to right. Let \(\mathcal T_n\) be the finite set of such trees with
\(n\) leaves. At a leaf carrying \(\mathsf{s}\xi_I\), place
\(\mathsf{s}\theta_I\). At each internal vertex use
\[
 \mathcal M(\mathsf{s}f,\mathsf{s}g)=(-1)^{|f|}
                                      \mathsf{s}(f\star g).
\]
On every internal edge place \(-\mathsf{s}h\mathsf{s}^{-1}\).
At the root place \(\mathsf{s}p\mathsf{s}^{-1}\) for an operation,
or \(-\mathsf{s}h\mathsf{s}^{-1}\) for a comparison component.
Compose these maps from leaves to root using the Koszul tensor rule.
Denote the resulting maps by \(B_T\) and \(F_T\), respectively.

\begin{theorem}\label{thm:co-all-inputs}
For every \(n\ge2\) and every sequence of exterior basis inputs,
the ordered transferred operations and comparison components are
\[
 b_n=\sum_{T\in\mathcal T_n}B_T,\qquad
 F_n=\sum_{T\in\mathcal T_n}F_T.
\]
Together with \eqref{eq:co-symbol-star}, these are finite formulas
using only derivatives of the potential, divided differences,
coefficient multiplication, and exterior signs.
They include unit inputs and arbitrary products of primitive generators.
They are coefficient-linear and continuous at each fixed arity in
the formal topology.
\end{theorem}

\begin{proof}
A tree has a unique root split into a left subtree with \(j\) leaves
and a right subtree with \(n-j\) leaves. Expanding the finite tree
sum by this split gives exactly the recurrence in
Section~\ref{sec:ot-contraction}. Induction from the one-leaf inclusion
therefore identifies every term, including its suspension sign.
This also identifies the unit-containing values with the strict-unit
identities proved there. The same contraction proof establishes
\(b^2=0\) and the full comparison equation.

Expanding each vertex by \eqref{eq:co-symbol-star} and each edge
by the displayed homotopy removes every unspecified operation from
the formula. There are finitely many trees, exterior symbols, and
contractions in each arity. Products and evaluations are continuous.
Each divided difference decreases formal order by at most one,
so every resulting fixed-arity map is continuous.
\end{proof}

The left ideal \(\mathcal J\) simplifies the root further.
Since \(F_m\) takes values in \(\mathsf{s}\mathcal J\) for \(m\ge2\),
\[
 b_n=\mathsf{s}p\mathsf{s}^{-1}
              \mathcal M(F_{n-1}\otimes F_1),\qquad n\ge3.
\]
Every other root split has its right factor in \(\mathcal J\)
and projects to zero. Internal subtrees can still have both
inputs composite. Their contractions account for products such
as \eqref{eq:co-first-composite}.

\section{Counting contractions and bounding arity}
\label{sec:co-degree}

Multiplication has two kinds of nonzero scalar contraction.
A pair of \(H\)'s contributes \(c_{ij}\), which is linear in
the potential. A pair \(\beta_i,H_i\) contributes \(1\).
Every homotopy introduces one \(\beta\) and one divided difference.
These three facts constrain each coefficient of the full formula.

For inputs \(\xi_{I_1},\ldots,\xi_{I_n}\), put
\(N=\sum_a|I_a|\). If a polynomial has degree at most \(m\),
degree means total degree in the variables
\(x,y\), or equivalently \(r,y\). Numerical coefficients of the
polynomial have degree zero in this convention.

\begin{theorem}\label{thm:co-degree}
For \(n\ge2\), the coefficient of \(\mathsf{s}\xi_K\) in
\(b_n(\mathsf{s}\xi_{I_1},\ldots,\mathsf{s}\xi_{I_n})\)
is homogeneous of degree
\[
 q=\frac{N-n+2-|K|}{2}
\]
in the coefficients of the potential. It is zero unless \(q\)
is a nonnegative integer. If \(\deg W\le m\), with \(m\ge2\),
its degree in \(y\) is at most
\[
 q(m-2)-(n-2).
\]
The coefficient of \(\mathsf{s}H_K\beta_J\) in \(F_n\) has
potential degree
\[
 q=\frac{N-n+1+|J|-|K|}{2}
\]
and degree in \(r,y\) at most \(q(m-2)-(n-1)\).
A negative upper bound forces the coefficient to vanish.
The statements hold over every commutative rational algebra.
\end{theorem}

\begin{proof}
Fix a term in the fully expanded tree formula. A term contributing
to \(b_n\) contains \(n-2\) homotopies. It therefore introduces
\(n-2\) symbols \(\beta\). Projection removes the term unless
all these symbols have contracted with \(H\)'s.
If \(q\) pairs of \(H\)'s also contract, the number of surviving
\(H\)'s is
\[
 |K|=N-(n-2)-2q.
\]
Exactly those \(q\) pairs contribute coefficients linear in \(W\).
Divided differences and evaluation do not change their degree under
the replacement \(W\mapsto tW\), for an even scalar parameter \(t\).
Every term contributing to the same output thus has potential degree \(q\).

For \(F_n\), there are \(n-1\) homotopies, and \(|J|\) of
their introduced symbols survive. Hence
\[
 |K|=N-(n-1-|J|)-2q,
\]
which gives its asserted potential degree.

When \(\deg W\le m\), every \(c_{ij}\) has total degree at most
\(m-2\). Multiplication adds degrees. A nonzero divided difference
decreases total degree by one, and evaluation does not increase it.
Following each tree from its leaves gives the two degree bounds.
The proof uses coefficient identities, not division by elements of \(k\).
It therefore also applies when \(k\) has zero divisors.
\end{proof}

The same counts apply to antisymmetrized representatives: a contraction
still removes two \(H\)'s and contributes one coefficient linear in \(W\).
Reversing the order of the normal homotopy does not change the counts.

\begin{corollary}\label{cor:co-finite-arity}
Set \(a=m-d(m-2)\). If \(a>0\), then on every input
\[
 b_n=0\quad\text{for }n>\frac{2m}{a},\qquad
 F_n=0\quad\text{for }n>\frac{m+d(m-2)}{a}.
\]
For every potential of degree at most three in two variables,
\(b_n=0\) for \(n\ge7\) and \(F_n=0\) for \(n\ge6\).
For one variable of degree at most \(m\), the bounds are
\(b_n=0\) for \(n>m\) and \(F_n=0\) for \(n\ge m\).
\end{corollary}

\begin{proof}
For basis inputs, \(N\le dn\) and \(|K|\ge0\).
The first degree bound can be nonnegative only if
\[
 2(n-2)\le\bigl((d-1)n+2\bigr)(m-2).
\]
Rearranging gives \(an\le2m\).
For the comparison, use also \(|J|\le d\).
Its degree bound requires
\[
 2(n-1)\le\bigl((d-1)n+d+1\bigr)(m-2),
\]
or \(an\le m+d(m-2)\).
Coefficient-linearity extends the result to all inputs.
Substituting \((d,m)=(2,3)\) or \(d=1\) gives the final assertions.
\end{proof}

For a quadratic potential these inequalities recover a strict
differential Clifford model. When \(a\le0\), the inequalities
give no uniform bound in arity. They do not assert that infinitely
many operations are nonzero.

\section{The highest operation for a binary cubic}
\label{sec:co-binary-cubic}

Take
\[
 W(x_1,x_2)=ax_1^3+bx_1^2x_2+cx_1x_2^2+dx_2^3
               +ex_1^2+fx_1x_2+gx_2^2+\ell x_1+m x_2+v_0,
\]
where all coefficients belong to \(k\). Write
\(r=x_1-y\), \(s=x_2-z\), and put
\begin{gather*}
 u=3ay+bz+e,\qquad v=cy+3dz+g,\qquad w=2by+2cz+f,
 \\
 P=bd-c^2,\qquad Q=bc-ad.
\end{gather*}
The ordered primitives have relations
\[
 H_1^2=-ar-bs-u,\qquad H_2^2=-ds-v,\qquad
 \{H_1,H_2\}=-cs-w.
\]
The mixed relations with \(\beta_1,\beta_2\) remain those of
Section~\ref{sec:co-composite-example}.

\begin{proposition}\label{prop:co-sixth}
For \(\chi=\mathsf{s}(\xi_1\xi_2)\), the only possible
nonzero sixth operation on nonunit basis inputs is
\[
 b_6(\chi,\chi,\chi,\chi,\chi,\chi)=ac(bd-c^2)e.
\]
The highest comparison component is
\[
 F_5(\chi,\chi,\chi,\chi,\chi)
       =-ac(bd-c^2)\mathsf{s}(\beta_1\beta_2).
\]
Both arity bounds of Corollary~\ref{cor:co-finite-arity} are sharp.
\end{proposition}

\begin{proof}
Let \(K_j=\mathsf{s}^{-1}F_j(\chi,\ldots,\chi)\), so
\(K_1=H_1H_2\). Since \(\chi\) is odd, the recurrence gives
\[
 K_n=-h\sum_{j=1}^{n-1}(-1)^{j+1}K_jK_{n-j}.
\]
Applying the displayed relations and the two divided differences yields
\begin{align*}
 K_2={}&a(ds+v)\beta_1+(bd s+bv+ud)\beta_2
                         +cH_1H_2\beta_2,\\
 K_3={}&P H_1\beta_2+acH_2\beta_1+QH_2\beta_2,\\
 K_4={}&-aP\beta_1+d(2ac-b^2)\beta_2
                         +ac^2H_2\beta_1\beta_2,\\
 K_5={}&-acP\beta_1\beta_2.
\end{align*}
The reductions can be checked without expanding the full matrices.
For \(K_2\), use
\((H_1H_2)^2=-H_1^2H_2^2+\{H_1,H_2\}H_1H_2\),
then remove the constant term in \(r\), followed by that in \(s\).
For \(K_3\), the sum before applying \(-h\) is
\[
 (Ps+bv+ud-cw)H_1+(ac r+Qs+cu-av)H_2.
\]
Its two divided differences give the displayed value, with the
minus sign of \(h\) on each odd \(H_i\).

For the next two reductions the sums before the homotopy are
\(K_1K_3-K_2^2+K_3K_1\) and
\(K_1K_4-K_2K_3+K_3K_2-K_4K_1\).
Their nonzero homotopy values are respectively
\[
 aP\beta_1-d(2ac-b^2)\beta_2-ac^2H_2\beta_1\beta_2,
 \qquad acP\beta_1\beta_2.
\]
These equalities follow by moving each \(\beta_i\) to the right
with \(\beta_iH_j=\delta_{ij}-H_j\beta_i\), and using
\(H_2H_1=-H_1H_2-cs-w\).
Terms already containing \(\beta_1\) have zero homotopy.
On terms containing \(\beta_2\), only their \(r\)-difference
contributes. On terms without a \(\beta\), take the \(r\)-difference
and then the \(s\)-difference after setting \(r=0\).
This specifies each scalar reduction and proves the formulas for
\(K_4,K_5\).

At the root of the sixth operation, only \(K_5K_1\) survives
projection. The scalar part of
\(\beta_1\beta_2H_1H_2\) is \(-1\).
The root suspension sign is positive, giving \(acP e\).
For any other nonunit sixth input word, \(N\le11\).
The degree bound in Theorem~\ref{thm:co-degree} requires
\(N\ge12+|K|\) for a nonzero output. Hence only the six
composite inputs and scalar output can occur.
The choice \(W=x_1^3+x_1x_2^2\) gives \(K_5=\beta_1\beta_2\)
and \(b_6(\chi^6)=-e\), proving sharpness.
\end{proof}

The coefficient \(ac(bd-c^2)\) belongs to this specified ordered
contraction. It is not invariant under changes of coordinates.
For \(W=x_1^3+x_2^3\) it is zero. The invertible substitution
\(x_1=s+t\), \(x_2=s-t\) gives \(W=2s^3+6st^2\),
whose coefficient is \(-432\).
Thus a comparison after a linear change of the field coordinates
can require higher components on the transferred algebra.
A displayed sixth multiplication alone cannot decide formality.
Its transformation law, and any cyclic structure that the comparison
must preserve, remain part of that question.
